% from main.tex
%%%%%%%%%%  scan idx 534  |  folio 521  %%%%%%%%%%
pour tout $x\in X(n)$, et on a
\begin{equation}\tag{1.1.6}
\ell_x=\ell'_x+\ell''_x\qquad (x\in X(n)).
\end{equation}
Alors l'assertion (1.1.3) pour tous les membres de (1.1.5) découle immédiatement de l'hypothèse et de l'équation
\begin{equation}\tag{1.1.7}
\cl_\bullet(G)=\cl_\bullet(G')+\cl_\bullet(G'').
\end{equation}

On sait que $K_\bullet$ est un foncteur covariant sur la catégorie des schémas localement noethériens par rapport aux morphismes propres. La proposition suivante affirme que notre filtration est fonctorielle :

\textbf{Proposition 1.1.8.} Soit $f:X'\to X$ un morphisme propre de schémas localement noethériens et soit
\[
f_*:K_\bullet(X')\longrightarrow K_\bullet(X)
\]
le morphisme induit par $f$. Alors
\begin{equation}\tag{1.1.9}
f_*(\Filt_n(X'))\subset\Filt_n(X)
\end{equation}
pour tout entier $n$.

\emph{Démonstration :} Si $F$ est un $\mathcal O_{X'}$-module cohérent, alors on a par définition (IV 2.11) :
\[
f_*(\cl_\bullet(F))=\sum_{i\in\mathbb N}(-1)^i\cl_\bullet(R^if_*(F)).
\]
Or, les $\mathcal O_X$-modules $R^if_*(F)$ sont concentrés sur $f(\Supp(F))$ et d'après (EGA IV 5.4.1 (ii)) on a
\[
\dim(\Supp(F))\geq\dim(f(\Supp(F))),
\]
d'où l'assertion.

\subsection*{1.2}
